\documentclass[11pt,a4paper]{article}

\usepackage[utf8]{inputenc}
\usepackage[T1]{fontenc}
\usepackage{amsmath,amssymb,amsfonts,amsthm}
\usepackage{geometry}
\usepackage{hyperref}
\usepackage{cite}
\usepackage{microtype}

\geometry{margin=1in}

\hypersetup{
    colorlinks=true,
    linkcolor=blue,
    filecolor=magenta,
    urlcolor=cyan,
    citecolor=blue,
}

\usepackage{perpage}
\MakePerPage{footnote} % reset footnote numbering on every page

\renewcommand*\footnoterule{} % no line before footnotes

\usepackage[symbol,bottom,hang,flushmargin]{footmisc}

\setlength{\skip\footins}{2em plus1em} % space between the text body and the footnotes
\setlength{\footnotesep}{\baselineskip} % space between footnotes

\newcommand\fex{\hspace{.2ex}} % one fifth of the ex
\newcommand\nfex{\hspace{-0.2ex}} % negative one fifth
\newcommand\thex{\hspace{.1ex}} % one tenth of the ex
\newcommand\nthex{\hspace{-0.1ex}} % negative one tenth

\newcommand\cin{\nfex\in\nfex}
\newcommand\cto{\nfex\to\nfex}

% a custom style with no trailing period
\newtheoremstyle{noperiod}
  {\topsep}   % space above
  {\topsep}   % space below
  {\itshape}  % body font
  {}          % indent amount
  {\bfseries} % theorem head font
  {}          % punctuation after theorem head (leave empty for no period)
  {.5em}      % space after theorem head
  {}          % theorem head spec

\theoremstyle{definition}
\newtheorem{definition}{Definition}[section]
\theoremstyle{noperiod}
\newtheorem{axiom}{Axiom}[section]
\newtheorem{theorem}{Theorem}[section]
\newtheorem{lemma}[theorem]{Lemma}

\title{\textbf{Limitless Summation}}

\author{Vadique, with the assistance of electronic intelligence}
\date{\today}

\begin{document}

\maketitle

\begin{abstract}
In traditional analysis, the concept of an~infinite series is conflated with its metric evaluation,
typically defined via the limit of its partial sums. This paper presents an~alternative framework
that decouples the infinite sum as a~formal algebraic object from its final numerical value.
By treating the infinite series as a~primary, standalone structural entity, the foundational paradigm
shifts from metric distance measurements defined by the theory of~limits to pure algebraic syntax.
Within this framework, a~consistent arithmetic for transfinite elements is introduced,
effectively bypassing historical paradoxes and singularities.
\end{abstract}

\thispagestyle{empty}

\section{Introduction}

\subsection{The Classical Conflation: Syntax vs. Metric Semantics}

In textbook mathematical analysis, an~expression representing an~infinite addition is not imagined as an~independent algebraic entity, precluded from existing outside its limit value.
Since standard arithmetic defines addition as a~binary operation, an~infinite sequence of~additions cannot be computed through step-by-step evaluation.
To assign meaning to an~infinite series $\sum a_n$, Cauchy and his successors bound the structural string to a~metric valuation tool, namely the limit of a~sequence of finite ``partial'' sums.
If this metric limit fails to settle upon a~single finite value, the infinite sum is declared ``divergent'' and dismissed as mathematically intractable.
This creates a~fundamental conflation between the syntax of an~infinite expression and its semantic measurement on a~real line.

\subsection{Historical Evolution of Infinite Summation}

The treatment of infinite addition has evolved through three major paradigms before its nineteenth-century crystallization.

\subsubsection{Antiquity and the Method of Exhaustion}

In ancient Greek mathematics, actual infinity was structurally taboo, driven by the logical paradoxes of Zeno of Elea. When Archimedes of Syracuse computed the area under a~parabolic segment (c. 250 BCE) via the geometric progression $1 + \frac{1}{4} + \frac{1}{16} + \dots = \frac{4}{3}$, he did not define an~infinite sum. Instead, he utilized the Method of Exhaustion—a double-contradiction proof demonstrating that the geometric area could neither be strictly greater than nor strictly less than $\frac{4}{3}$. Infinity was bypassed by maintaining strictly finite, arbitrarily large metric bounds.

\subsubsection{The Merton School, Oresme, and Structural Groupings}

In the 14th century, the study of kinematics forced Scholastic calculators to analyze continuous motion. The Oxford Calculators at Merton College attempted to track total distance as an~aggregate of changing velocities. Concurrently, Nicole Oresme (c. 1350) provided the first formal proof of the divergence of the Harmonic series. Crucially, Oresme did not use limits; he utilized a~pure structural grouping technique:
\begin{equation}
1 + \frac{1}{2} + \biggl(\frac{1}{3} + \frac{1}{4}\biggr) + \biggl(\frac{1}{5} + \dots + \frac{1}{8}\biggr) > 1 + \frac{1}{2} + \frac{1}{2} + \frac{1}{2} + \dots
\end{equation}
The expression was treated as an~operational recipe generating infinite discrete blocks of mass, proving properties of the string through its global arrangement rather than localized values.

\subsubsection{The Eulerian Era: Infinite Polynomials and Formal Algebra}

Following the development of calculus by Newton and Leibniz, the 17th and 18th centuries treated infinite series using formal algebraic grammar. Mathematicians such as Leonhard Euler and the Bernoulli brothers explicitly operated on infinite series as "infinite polynomials." They did not restrict operations to domains of convergence. Euler famously solved the Basel Problem ($\sum \frac{1}{n^2} = \frac{\pi^2}{6}$) by factoring the infinite Taylor expansion of $\sin(x)$ into an~infinite product of its roots, treating the transfinite string exactly like a~finite algebraic expression.

\subsection{The 19th-Century Metric Takeover}

This unconstrained algebraic fluency eventually caused systemic paradoxes, most famously illustrated by Grandi's alternating series $1 - 1 + 1 - 1 \dots$, which could be grouped to equal $0$, $1$, or $\frac{1}{2}$ depending on shifting parentheses. Panicked by these contradictions, 19th-century analysis enacted a~complete metric takeover.

Augustin-Louis Cauchy, in his famous \textit{Cours d'analyse} (1821), redefined the infinite series by tethering it to a~finite value, calling a~series convergent (\textit{série convergente}) if its partial sums approached a~fixed limit indefinitely. This dynamic view can be formulated topologically as
\begin{equation}
\textstyle\sum_{k=1}^{\infty} a_k = s \iff n \nfex\uparrow \:\: \Rightarrow \Bigl( \textstyle\sum_{k=1}^{n} a_k - s \Bigr) \nfex\cin \forall\fex \mathsf{U}
\end{equation}
where ${\forall\fex \mathsf{U}}$ is any~(\textit{indéfiniment}) open neighborhood of~$0$.

Later, during his university lectures in Berlin (c. 1859–1861), Karl Weierstrass arithmetized this definition. Instead of publishing a~formal standalone paper, Weierstrass introduced his parameters through oral pedagogical networks. He inverted Cauchy's dynamic phrase \textit{``des valeurs de n toujours croissantes''} into a~static, functional constraint requiring the determination of an~inverse threshold function $N(\varepsilon)$ such that for all $n > N$:
\begin{equation}
\left|\, \sum_{k=1}^{n} a_k - s \,\right| < \varepsilon
\end{equation}
This framework successfully policed boundaries but did so by entirely eliminating formal algebraic operator fluency at the transfinite edge.

\subsection{Motivation: The Inversion Problem and Pedagogical Intractability}

The modern pedagogical reliance on the Cauchy-Weierstrass framework introduces a~severe analytical friction point. To prove convergence under the standard definition, one must structurally invert a~discrete function mapping $\mathbb{N} \cto \mathbb{R}$. However, for the vast majority of non-trivial or transcendental series, finding an~exact inverse function $N(\varepsilon)$ is analytically intractable. Consequently, modern mathematical practice relies on loose metric upper bounds, integral approximations, or pre-calculated remainder coefficients ($R_n$) to satisfy the definition after the fact.

This paper proposes that this inversion process is an~artificial requirement of metric distance tracking. By evaluating the infinite series directly at a~completed transfinite boundary ($\aleph$), we eliminate the need for metric inversion functions, resolving classical divergences into stable structural coordinates.

\section{Philosophical and Topological Framework}

\subsection{The Concept of the Completed Path Continuum}

We reject the notion that an~infinite sum is an~endless, incomplete chronological addition process. Instead, we model an~infinite series as a~completed path continuum. Just as modern set theory accepts a~countable infinity as a~completed set ($\mathbb{N}$), we accept the operations over that set as a~completed transfinite object.

\subsection{Defining the Infinite Sum as a~Primary Standalone Mapping Object}

Let $\mathbb{F}$ represent a~standard base field ($\mathbb{Q}$, $\mathbb{R}$, or $\mathbb{C}$), and let $\mathbb{Z}$ represent the index domain. An~infinite series is defined as a~functional discrete mapping $f: \mathbb{Z} \to \mathbb{F}$. We define the \textit{Limitless Summation Operator}, denoted by $\mathcal{S}$, as a~primary, standalone structural operator:
\begin{equation}
\mathcal{S}: (\mathbb{Z} \to \mathbb{F}) \longrightarrow \mathbb{W}_{\mathbb{F}}
\end{equation}
where $\mathbb{W}_{\mathbb{F}}$ is a~generalized non-Archimedean field-wheel extension. The sum $\mathcal{S}[f]$ evaluates the continuous mass of the function up to the transfinite boundary coordinate $\aleph$. The output is a~singular, fixed entry within $\mathbb{W}_{\mathbb{F}}$, independent of any real-line metric measurement.

\subsection{Definition of the Structural Transfinite Logarithm as a~Density Counter}

To govern sub-sampled collections of terms without changing the underlying variable matrix, we introduce a~novel structural operator.
\begin{definition}
The \textit{Structural Transfinite Logarithm}, denoted by $\log(\aleph)$, is defined as the cardinal count of the sparse subset of indices $n \le \aleph$ that satisfy the Diophantine approximation threshold dictated by Dirichlet’s Approximation Theorem for an~irrational constant.
\end{definition}
By definition of Diophantine sparsity, $\log(\aleph)$ acts as a~structural density coefficient mapping the relative distance between elements along the linear continuum, satisfying the scale dominance identity:
\begin{equation}
\frac{\log(\aleph)}{\aleph^k} \equiv 0 \quad \forall k \ge 1
\end{equation}

\section{The Algebraic Specification of the Extended Field-Wheel $\mathbb{W}_{\mathbb{F}}$}

\subsection{Closure and Domain Formalization}

To guarantee complete algebraic closure under transfinite operations, the output space of the summation operator is extended beyond the base scalar field $\mathbb{F}$.
We define the Generalized Field-Wheel Extension $\mathbb{W}_{\mathbb{F}}$ as:
\begin{equation}
\mathbb{W}_{\mathbb{F}} = \mathbb{F} \cup \{ \log(\aleph), \aleph, 2^\aleph \} \cup \{ \%, \infty \}
\end{equation}
where $\%$ represents the nullary un-evaluatable zero (``nullity'' $\frac{0}{0}$), and $\infty$ represents the absolute infinity.

\subsection{The Primary Structural Axioms}

The operational grammar of $\mathbb{W}_{\mathbb{F}}$ relaxes classical field cancellation rules to allow division by zero and boundary stabilization, governed by three primary axioms.

\begin{axiom}[Linear Shift Invariance]
\begin{equation}
\aleph \pm n = \aleph \quad \forall n \in \mathbb{Z}
\end{equation}
Stepping backward or forward by a~finite integer step does not alter the coordinate of the completed path at the transfinite boundary.
\end{axiom}

\begin{axiom}[Projective Inversion and Loop Closure]
\begin{equation}
\aleph \cdot (-1) = \aleph
\end{equation}
Reversing the structural direction of a~completed transfinite path yields an~isomorphic, identical object. This topologically closes the transfinite edge into an~unsigned projective pole.
\end{axiom}

\begin{axiom}[Global Exponentiation Idempotency]
\begin{equation}
x^0 = 1 \quad \forall x \in \mathbb{W}_{\mathbb{F}}
\end{equation}
This identity holds unconditionally for all finite and transfinite quantities, including zero $(0^0 = 1)$.
\end{axiom}

\subsection{The Variable Zero and Path Preservation}

To bypass Bishop Berkeley’s historical ``ghosts of departed quantities'' critique of early calculus, we establish a~rigid distinction between a~dead constant zero and a~dynamic variable zero.

\begin{axiom}
Let $x$ represent a~continuous algebraic variable mapping through a~boundary threshold. The identity expression satisfies:
\begin{equation}
\frac{x}{x} = 1
\end{equation}
globally, including at the evaluation point $x = 0$.
\end{axiom}
Because the variable carries its structural path history with it, fractions must be simplified algebraically prior to static coordinate evaluation, neutralizing singularities and removing holes from the functional continuum.

\subsection{Working with Transfinite Fractions}

Arithmetic operations within the field-wheel extension follow a~strict hierarchical calculus:
\begin{enumerate}
\item \textbf{The Principle of Finite-Stage Priority:} All algebraic factorizations, cancellations, and simplifications must be fully applied at the finite variable stage ($n$) before the substitution of the transfinite coordinate $\aleph$.
\item \textbf{The Non-Cancellation of Free Ratios vs. Bound Variables:} Unlinked transfinite elements do not satisfy field cancellation. For example, the free transfinite mass ratio of~unlinked layers satisfies $\frac{\aleph}{\aleph} \neq 1$ and represents independent unreduced transfinite rational expression. However, if the numerator and denominator are bound to the \textit{exact same variable within a~single monomial mapping} originating from the finite stage, cancellation remains valid upon transfinite substitution, for example, $\aleph \text{ times } \bigl[ \frac{1}{n^2} \bigr]_{n=\aleph} \nfex = \frac{\aleph}{\aleph^2} = \frac{1}{\aleph}$.
\item \textbf{Non-Invertibility of Scale Dominance:} The scale dominance identity $\frac{\aleph}{2^\aleph} = 0$ is non-invertible due to the absorbing nature of wheel zeros. It does not imply $2^\aleph \cdot 0 = \aleph$, nor does it imply $2^\aleph = \frac{\aleph}{0}$. Under wheel rules, division by zero is defined via the projective boundary ($\frac{\aleph}{0} = \aleph$).
\item \textbf{Constant Scale Absorption:} For any finite non-zero scalar~${k \cin \mathbb{F}}$, the counting path absorbs density scaling: $k\aleph = \aleph$.
\end{enumerate}

\section{Foundational Structural Evaluations}

\subsection{The Harmonic Series: Boundary Tail Accumulation}

We deploy our transfinite calculus to evaluate the classical Harmonic series:
\begin{equation}
\mathcal{S}\left[\frac{1}{n}\right] = 1 + \frac{1}{2} + \frac{1}{3} + \frac{1}{4} + \dots
\end{equation}
Evaluating the series from ``that end'' requires stepping directly onto the completed transfinite coordinate $\aleph$ and looking backward across the infinite tail:
\begin{equation}
\text{Tail Mass} = \frac{1}{\aleph} + \frac{1}{\aleph-1} + \frac{1}{\aleph-2} + \dots + \frac{1}{\aleph-k}
\end{equation}
Applying the linear shift invariance axiom ($\aleph - c = \aleph$), every denominator in the infinite tail collapses identically into the stable boundary coordinate, creating a~completely uniform field:\begin{equation}
\text{Tail Mass} = \frac{1}{\aleph} + \frac{1}{\aleph} + \frac{1}{\aleph} + \dots + \frac{1}{\aleph}
\end{equation}
Because completing this path requires exactly $\aleph$ steps, the term $\frac{1}{\aleph}$ is accumulated precisely $\aleph$ times. This yields the total structural mass of the series:
\begin{equation}
\mathcal{S}\left[\frac{1}{n}\right] = \frac{\aleph}{\aleph}
\end{equation}

The Harmonic series does not blow up into an~unmeasurable void, it maps cleanly to a~stable, unreduced transfinite rational coordinate within $\mathbb{W}_{\mathbb{F}}$.

\subsection{The Euler--Basel Series: Boundary Mass Decay}

We apply the backward tail technique to the reciprocal squares series:
\begin{equation}
\mathcal{S}\left[\frac{1}{n^2}\right] = 1 + \frac{1}{4} + \frac{1}{9} + \frac{1}{16} + \dots
\end{equation}
Transitioning to the transfinite coordinate $\aleph$ and executing the backward path expansion yields:
\begin{equation}
\text{Tail Mass} = \frac{1}{\aleph^2} + \frac{1}{(\aleph-1)^2} + \frac{1}{(\aleph-2)^2} + \dots \implies \sum_{k=1}^{\aleph} \frac{1}{\aleph^2}
\end{equation}

Accumulating this uniform quadratic field across $\aleph$ transfinite steps isolates the total mass profile at the edge of the universe:
\begin{equation}
\text{Total Boundary Mass} = \aleph \cdot \left(\frac{1}{\aleph^2}\right) = \frac{\aleph}{\aleph^2}
\end{equation}
Because both terms in the ratio $\frac{\aleph}{\aleph^2}$ are bound to the identical, unmodified index variable~$n$ inherited from the finite stage, this constitutes a~valid structural identity cancellation, reducing the expression cleanly to
\begin{equation}
\frac{\aleph}{\aleph^2} \equiv \frac{1}{\aleph}
\end{equation}

Under scale dominance axioms, an~expression leaving the linear counter strictly in the denominator constitutes an~infinitesimal that resolves to a~structural zero ($\frac{1}{\aleph} \equiv 0$). The quadratic dampening completely neutralizes the transfinite boundary mass, allowing the series to shed all transfinite residue and collapse smoothly into a~standard finite value inside the base field $\mathbb{F}$ (measured classically as $\frac{\pi^2}{6}$).

\subsection{The Indeterminate Geometric Series: Algebraic Reciprocity}

We analyze the geometric progression $\Gamma(n)$ in an~indeterminate variable~$x$:
\begin{equation}
\Gamma(n) = 1 + x + x^2 + x^3 + x^4 + \dots + x^n
\end{equation}
We establish two simultaneous structural identities for the finite sequence:
\begin{align}
\text{Identity }\mathbf{A}: \quad \Gamma(n) &= \Gamma(n-1) + x^n \\
\text{Identity }\mathbf{B}: \quad \Gamma(n) &= 1 + x\Gamma(n-1)\end{align}
Equating both expressions via substitution yields the finite relation:
\begin{equation}
x^n - 1 = (x-1)\Gamma(n-1)
\end{equation}

We now evaluate this system at the completed transfinite boundary $n = \aleph$. Utilizing Identity $\mathbf{A}$ at the boundary, our shift invariance axiom mandates that $\Gamma(\aleph) = \Gamma(\aleph-1)$. Substituting this identity directly into Identity $\mathbf{A}$ reduces the system to:
\begin{equation}
\mathbf{A}(\aleph): \quad \Gamma(\aleph) - \Gamma(\aleph) = x^\aleph \implies x^\aleph = 0
\end{equation}
This establishes a~strict structural constraint: for the geometric sum to maintain algebraic closure, the boundary term $x^\aleph$ must act as a~nilpotent-like element resolving to zero. Evaluating Identity $\mathbf{B}$ at the transfinite boundary under this constraint yields:
\begin{equation}
\mathbf{B}(\aleph): \quad 0 - 1 = (x-1)\Gamma(\aleph) \implies \Gamma(\aleph) = \frac{1}{1-x}
\end{equation}
The expression $\frac{1}{1-x}$ is derived as a~global algebraic identity independent of metric boundaries. For $x=2$, the sum $1+2+4+8+\dots = -1$ if and only if $2^\aleph = 0$ is true in the underlying number system. This conditional constraint maps perfectly to the $2$-adic integers ($\mathbb{Q}_2$), proving that Limitless Summation adaptively adjusts its grammar to the destination field.

\section{Advanced Evaluation of the $\sin^2(n)$ Family Series}

\subsection{Pathological Obstacles in Real Analysis}

The series $\sum \frac{1}{n^2 \sin^2(n)}$ constitutes a~famous pathological nightmare in classical limit theory. Because the integer argument is in radians, and the constant $\pi$ is irrational, the denominator periodically approaches zero when $n$ hits a~near-multiple of $\pi$. Classical limits are forced to track these chaotic metric spikes via Diophantine approximations, causing the standard series to violently diverge. We now demonstrate how our transfinite wheel regularizes this expression.

\subsection{Derivation I: The Geometric Boundary Condition and the Variable Constant $c^2$}

Because $\pi \notin \mathbb{Q}$, it is an~absolute number-theoretic identity that for any integer index $n \in \mathbb{Z}$, an~integer can never perfectly equal a~multiple of radians matching the zero-crossings of the wave:
\begin{equation}
n \neq k\pi \quad \forall n, k \in \mathbb{Z} \implies \sin^2(n) \neq 0 \quad \forall n \in \mathbb{N}
\end{equation}
When evaluated at the completed transfinite coordinate $\aleph$, the index ceases to be a~wandering variable. Under the identity $\aleph = \aleph - 1$, the index is frozen into a~static boundary state, smoothing out the chaotic oscillations. Therefore, the boundary term $\sin^2(\aleph)$ must stabilize into a~non-zero transfinite constant, which we denote as $c^2$.
Accumulating $\aleph$ copies of these quadratically dampened terms reveals that the total mass scales as
\begin{equation}
\text{Total Mass} = \frac{\aleph}{\aleph^2 c^2}
\fex ,
\end{equation}
where the plain, unmodified index counters cancel to yield
\begin{equation}
\frac{\aleph}{\aleph^2 c^2} \equiv \frac{1}{\aleph c^2}
%%\fex .
\end{equation}

\subsection{Derivation II: Multi-Layer Infinite Series Decomposition}

To verify the structural mechanics of $c^2$, we decouple the denominator by expanding $\sin^2(n) = \frac{1}{2} - \frac{1}{2}\cos(2n)$ into its formal Taylor power series, evaluating it at the unlinked nested infinite boundaries $n = \aleph$ and $k = \aleph$:
\begin{equation}
\sin^2(n) = \sum_{k=1}^{\infty} (-1)^{k-1} \frac{2^{2k-1}}{(2k)!} n^{2k}
\end{equation}
Following the \textbf{Cantor-Bernstein-Schröder Theorem (Bernstein, 1898)}, the transfinite factorial acts as a~higher-tier exponential infinity, satisfying the structural bijection ${\aleph! = \aleph^\aleph = 2^\aleph}$.

At the boundary layer, the coefficients yield the free transfinite ratio $\frac{2^{2\aleph-1}}{(2\aleph)!} \implies \frac{\infty}{\infty}$, representing a~stable transfinite rational mass. Substituting our boundary coordinate $n = \aleph$ into the nested expansion creates a~direct interaction with the exponential denominators ($\aleph^\aleph = \infty$):
\begin{equation}
\sin^2(\aleph) \implies \left(\frac{\aleph}{\aleph}\right) \cdot (-1)^\aleph \cdot \left(\frac{\infty}{\infty^2}\right)
\end{equation}
Under our Axiom of Scale Dominance, the quadratic exponential tier ($\infty^2$) in the denominator completely outscales the linear transfinite components and the boundary oscillation $(-1)^\aleph$ in the numerator, forcing the free tail ratio to resolve to an~absolute structural zero:
\begin{equation}
\left(\frac{\aleph}{\aleph}\right) \cdot (-1)^\aleph \cdot \left(\frac{\infty}{\infty^2}\right) \equiv 0
\end{equation}
The volatile, oscillating tail terms of the sine expansion are crushed to zero by the factorial hierarchy, forcing $\sin^2(\aleph)$ to stabilize into a~non-zero transfinite constant.

\subsection{Derivation III: The Finitist Dirichlet Objection and the Peak Threshold}

A classical objector applying Archimedean metric logic states that by Dirichlet's Approximation Theorem, there exist infinitely many integer spikes where $a_n > \frac{1}{16}$. They claim that a~partial sum must eventually contain $16S$ such terms, forcing the sum to exceed any bound $S$, thereby proving divergence without invoking infinity.

However, by the laws of Diophantine approximation (\cite{khinchin1964}), the spacing between these extraordinary peaks grows at an~exponential rate (\cite{rockett1992}). To accumulate a~finite count of $16S$ peaks, the partial sum's index must expand exponentially to a~threshold of $n = 2^{16S}$. This displacement is not finite; it sits precisely at the completed counting continuum $\aleph$. Because the total count of these sparse peaks up to the boundary is defined by our Structural Logarithm operator, we set the required peak threshold equal to the structural density of the highway:
\begin{equation}
16S = \log(\aleph) \implies S = \log(\aleph)
\end{equation}
by transfinite constant absorption rules.

\subsection{Synthesis: Resolving the Boundary Constant $c^2$}We unify all three derivations to solve for our missing geometric constant. By setting the total mass from Derivation I equal to the structural value derived from Derivation III, we obtain the fundamental identity of the series:
\begin{equation}
\frac{1}{\aleph c^2} = \log(\aleph) \implies c^2 = \frac{1}{\aleph \log(\aleph)}
\end{equation}
The boundary constant of the sine wave is exactly balanced by the product of the linear continuum and its density counter.

\section{Comparative Analysis and Harmonization of Paradigms}

\subsection{The Resolution of the Finitist Dilemma}

The confrontation between the two frameworks is revealed to be a~semantic shift at the transfinite boundary rather than an~arithmetic error. To gather $16S$ terms greater than $\frac{1}{16}$, one must push the index matrix to $n = 2^{16S} = \aleph$. At this coordinate, the total accumulated mass of the peak subset evaluates precisely to the structural transfinite logarithm, $\log(\aleph)$.
\begin{itemize}
\item \textbf{The Classical Framework:} Lacks the algebraic grammar to process transfinite operators as stable destinations, conflates $\log(\aleph)$ with an~unmanageable divergent infinity ($\infty$), and marks the series as \textbf{divergent}.
\item \textbf{The Operator Framework:} Recognizes $\log(\aleph)$ as a~well-defined, structurally regularized non-Archimedean transfinite coordinate. The series converges structurally to the transfinite value $\log(\aleph)$.
\end{itemize}

\subsection{The Definiteness of Boundary Mass: Quadratic vs. Cubic Variations}

We test the predictive power of this identity by evaluating the cubic variation $\sum \frac{1}{n^3 \sin^2(n)}$.
Substituting our derived boundary constant $c^2 = [\aleph \log(\aleph)]^{-1}$ into the cubic mass template yields
\begin{equation}
\text{Total Cubic Mass} = \frac{\aleph}{\aleph^3 c^2} = \frac{\aleph}{\aleph^3 \cdot \frac{1}{\aleph \log(\aleph)}} = \frac{\aleph^2 \log(\aleph)}{\aleph^3}
\end{equation}
Because the $\aleph^2$ in the numerator and the $\aleph^3$ in the denominator are strictly bound to mono-variate transformations of the identical finite index $n$, they undergo valid identity reduction.
The~expression simplifies cleanly to
\begin{equation}
\frac{\aleph^2 \log(\aleph)}{\aleph^3} \equiv \frac{\log(\aleph)}{\aleph}
\end{equation}
Under our Tier 3 Scale Dominance axioms, the linear counter dominantly outscales the structural logarithm, forcing the entire transfinite boundary expression to resolve to a~clean, absolute zero:
\begin{equation}
\frac{\log(\aleph)}{\aleph} \equiv 0
\end{equation}

This zero-resolution proves that the extra power in the cubic denominator completely neutralizes the transfinite boundary layer. While the quadratic variation accumulates enough mass to settle on the transfinite coordinate $\log(\aleph)$, the cubic variation sheds all transfinite residue, collapsing cleanly into a~standard finite value within the base field $\mathbb{F}$, perfectly matching classical expectations.

\section{Concluding Remarks and Architectural Outlook}

This paper has demonstrated that by treating the infinite series as a~primary, standalone mapping operator into the closed field-wheel extension $\mathbb{W}_{\mathbb{F}}$, we restore algebraic fluency to transfinite boundaries. Far from causing contradictions, Limitless Summation regularizes classically divergent expressions into precise transfinite identities.

The regularization of transfinite summation mass achieved in this work forms the baseline framework for a~broader restructuring of mathematical analysis. By demonstrating that variables can safely carry their bound structural path identities through boundary states—such as our definition of the dynamic variable zero ($\frac{x}{x}=1$ at $x=0$)—we establish the operational grammar required to bypass limits entirely. Future work will demonstrate that the differential and integral apparatus of calculus can be executed as pure, algebraic operator grammar via the shift operator $E = \exp(d)$, reclaiming the elegant, structural fluency pioneered by Leibniz and Euler.

\begin{thebibliography}{99}

\bibitem{khinchin1964}
A.~Y. Khinchin.
\newblock {\em Continued Fractions}.
\newblock University of Chicago Press, 1964.

\bibitem{rockett1992}
A.~M. Rockett and P.~Szüsz.
\newblock {\em Continued Fractions}.
\newblock World Scientific, 1992.

\bibitem{carlstrom2004}
J.~Carlström.
\newblock Wheels—on division by zero.
\newblock {\em Mathematical Structures in Computer Science}, 14(1):143--179, 2004.

\bibitem{kock2006}
A.~Kock.
\newblock {\em Synthetic Differential Geometry}, volume 333.
\newblock Cambridge University Press, 2006.

\bibitem{robinson1966}
A.~Robinson.
\newblock {\em Non-standard Analysis}.
\newblock North-Holland Publishing Company, 1966.

\bibitem{alekseyev2011}
M.~A. Alekseyev.
\newblock On the convergence of the Flint Hills series.
\newblock {\em arXiv preprint arXiv:1104.5100}, 2011.

\bibitem{zeilberger2020}
D.~Zeilberger and W.~Zudilin.
\newblock An~$\varepsilon$-improvement to the irrationality measure of~$\pi$.
\newblock {\em Mathematical Notes}, 108(3):450--454, 2020.

\end{thebibliography}

\end{document}
